\section{\kcode{taskgraph} Construct}

\label{sec:taskgraph}
\index{taskgraph construct@\kcode{taskgraph} construct}
\index{constructs!taskgraph@\kcode{taskgraph}}
\index{clauses!graph_id@\kcode{graph_id}}
\index{clauses!replayable@\kcode{replayable}}
\index{saved modifier@\kcode{saved} modifier}
\index{firstprivate clause@\kcode{firstprivate} clause!save modifier@\kcode{saved} modifier}

A taskgraph record represents a collection of tasks and their dependences
produced by a recorded sequence of task-generating constructs encountered
during the execution of a \kcode{taskgraph} construct, optionally identified by
a graph ID specified in a \kcode{graph_id} clause. While executing a region
corresponding to a \kcode{taskgraph} construct, the implementation may record
encountered task-generating constructs that are designated as replayable into a
taskgraph record. When a \kcode{taskgraph} construct is encountered and a
taskgraph record associated with the construct exists that matches any
specified graph ID), the implementation will execute from the taskgraph record
instead (referred to as a replay execution). 

In general, the \kcode{replayable} clause determines whether a task-generating
construct is replayable or not. But if the \kcode{replayable} clause is not
present on a task-generating construct that is lexically enclosed in a
\kcode{taskgraph} construct and is encountered during execution of the
\kcode{taskgraph} region (so, not during execution of another task generated
from the \kcode{taskgraph} region), then that construct is still considered
replayable.

Replay executions can provide some performance savings, as an implementation
can avoid setup code for task-generating constructs and identifying task
dependences. When executing from a taskgraph record, the implementation will
typically use the current enclosing data environment of the \kcode{taskgraph}
construct for variables in a variable list or locator list. Exceptions to this
are:

\begin{itemize}
\item For the \kcode{depend} clause, the storage locations are recorded and
      reused for replay executions.

\item For variable or locator list items that are composed of a base variable and
      various other expressions (e.g., an array element or array section), the
        values of the expressions are recorded and reused for replay execution.

\item If the \kcode{saved} modifier is present, a recorded copy of the variable
      in the taskgraph record is used.
\end{itemize}

In the following example, a \kcode{taskgraph} construct is used to allow the
implementation to produce a simple taskgraph record for replay execution.
Inside the construct, there are 3 \kcode{task} constructs that generate tasks
\plc{A}, \plc{B}, \plc{C}, respectively, where task \plc{C} is dependent on
tasks \plc{A} and \plc{B}. The tasks use the shared variables \ucode{x},
\ucode{y}, \ucode{z}, \ucode{sums}, and \ucode{i} from the data environment
that encloses the \kcode{taskgraph} construct. Given that \ucode{x}, \ucode{y},
\ucode{z}, and \ucode{sums} have static storage, the implementation can assume
it always comes from the same locations in memory for any replay execution of
the taskgraph record.

\cexample[6.0]{taskgraph}{1}

\ffreeexample[6.0]{taskgraph}{1}

The next example is a slight modification of the first one. The variables
\ucode{x}, \ucode{y}, and \ucode{z} are still shared in the recorded
\kcode{task} constructs, but now no longer have static storage as they are
declared local to the \ucode{do_parallel_work} procedure. The implementation
must take care of recapturing the addresses of the non-static \ucode{x},
\ucode{y}, and \ucode{z} prior to a replay execution of the taskgraph record
rather that referring to the storage of \ucode{x}, \ucode{y}, and \ucode{z}
from when the taskgraph record was created (which would no longer exist in the
replay executions).

\cexample[6.0]{taskgraph}{2}

\ffreeexample[6.0]{taskgraph}{2}

The next example again modifies the first one. Now, variables \ucode{x},
\ucode{y}, and \ucode{z} are declared as pointers to arrays declared in the
same scope in a procedure \ucode{exec_taskgraph}, which is called from the
\kcode{parallel} region. Consequentially, they are treated as firstprivate in
the recorded \kcode{task} constructs of the taskgraph record. During a replay
execution, firstprivate copies of these variables must be initialized from the
variables in the enclosing data environment of the current \kcode{taskgraph}
region. So, just as with the preceding example, the implementation must
recapture the addresses of the \ucode{x}, \ucode{y}, and \ucode{z} pointers
prior to a replay execution of the taskgraph record so that the firstprivate
copies are initialized by reading the pointers at the correct storage locations.

\cexample[6.0]{taskgraph}{3}

\ffreeexample[6.0]{taskgraph}{3}

The next example is much like the first one, except that \ucode{x}, \ucode{y},
and \ucode{z} are now static lifetime variables that are declared in the
\kcode{taskgraph} construct itself. They are shared in the recorded
\kcode{task} constructs, and as in the first example the implementation can
simply refer to their static storage without requiring an address recapture for
each subsequent replay execution.

\cexample[6.0]{taskgraph}{4}

\ffreeexample[6.0]{taskgraph}{4}

In the next example, \ucode{x}, \ucode{y}, and \ucode{z} are pointers into
\ucode{NG$\times$N} arrays, and the OpenMP implementation is directed to
create up to \ucode{NG} taskgraph records for replay execution. The pointers
\ucode{x}, \ucode{y}, and \ucode{z} are expected to point to the same array
slices for all replay executions of a given taskgraph record. The
\kcode{graph_id} clause specifies that the scalar integer \ucode{input_index},
which can have a value of 0 up to \ucode{NG}-1, is to be used as the graph ID.
Additionally, since the pointers \ucode{x}, \ucode{y}, and \ucode{z} don't
change for replay executions of a given taskgraph record, the example makes
them firstprivate in the recorded \kcode{task} constructs, and the use of the
\kcode{saved} modifier directs the implementation to initialize the
firstprivate copy from a value saved with the recorded \kcode{task} construct
in the taskgraph record. Note that if the \kcode{parallel masked} construct was
hoisted outside the \ucode{do_parallel_work} procedure and instead enclosed the
call, the behavior would remain the same.

\cexample[6.0]{taskgraph}{5}

\ffreeexample[6.0]{taskgraph}{5}
